\documentclass[11pt,a4paper]{article}

\usepackage[margin=1in]{geometry}
\usepackage{amsmath,amssymb,amsthm,mathtools,bm}
\usepackage{graphicx}
\usepackage{hyperref}
\usepackage{enumitem}
\usepackage{booktabs}
\usepackage{xcolor}
\usepackage{microtype}
\usepackage{mathrsfs}

\hypersetup{
  colorlinks=true,
  linkcolor=blue!60!black,
  citecolor=blue!60!black,
  urlcolor=blue!60!black
}

\title{\textbf{Phase Mirror: Prime–Indexed Phase--Dissonance}\\
Functionals for Software Governance}
\author{Citizen Gardens -- The Foundation of Multiplicity\\
The Prime Materia Commons\\
\texttt{info@citizengardens.org}}
\date{May 2026}

\newtheorem{theorem}{Theorem}[section]
\newtheorem{corollary}[theorem]{Corollary}
\newtheorem{proposition}[theorem]{Proposition}
\newtheorem{remark}[theorem]{Remark}
\newtheorem{definition}[theorem]{Definition}
\newtheorem{lemma}[theorem]{Lemma}

\newcommand{\Sig}{\mathrm{Sig}}
\newcommand{\Axes}{\mathrm{Axes}}
\newcommand{\globalSig}{\Sigma}
\newcommand{\weak}{\sim}
\newcommand{\strong}{\approx}
\newcommand{\Nat}{\mathbb{N}}
\newcommand{\Int}{\mathbb{Z}}
\newcommand{\Hilb}{\mathcal{H}}
\newcommand{\Tensor}{\mathcal{T}}
\newcommand{\Pow}{\mathcal{P}}
\newcommand{\op}{\mathrm{op}}
\newcommand{\PETC}{\mathrm{PETC}}
\newcommand{\PWEH}{\mathrm{PWEH}}
\newcommand{\ZRSD}{\mathrm{ZRSD}}
\newcommand{\PIRTM}{\mathrm{PIRTM}}
\newcommand{\diag}{\mathrm{diag}}
\newcommand{\spr}{\rho}
\newcommand{\Id}{I}

\begin{document}
\maketitle

\begin{abstract}
We present a prime-indexed recursive framework for tensor-structured mathematical objects in which constructive generation and contractive stabilization are treated as dual aspects of a single formal system. The constructive layer is supplied by the Prime-Encoded Tensor Constructor (\PETC), whose types are finite axis-signature lists over the free abelian group on a set of generators. The stabilizing layer is supplied by the Universal Multiplicity Constant \(\Lambda_m\), obtained as a scalar reduction of a sealed two-layer operator \(\Lambda_m^{\op}(t)\). Within this architecture we prove three core results: (i) reshape completeness for \PETC, namely that the equivalence relation generated by \(\{\mathrm{Permute},\mathrm{Merge},\mathrm{Split}\}\) is exactly equality of global signature; (ii) Banach contractivity of the prime-recursive meta-map under a finite prime truncation and a contractive parameter bound \(k<1\); and (iii) reachability of any admissible fixed point by suitable choice of driving term. We further situate commitment encodings and witness bytecode as a governance substrate for proof-carrying tensor transformations. Simulation-oriented components, including \ZRSD\ and enhanced tunneling dynamics, are stated as an implementation program rather than as completed empirical proof. The result is a unified but epistemically disciplined framework: \PETC\ supplies lawful construction, \(\Lambda_m\) supplies lawful contraction, and the fixed-point dynamics provide a recursive substrate for tensor-structured mathematical existence.
\end{abstract}
\newpage    
\tableofcontents
\newpage
\section{Introduction}

A long-standing foundational ambition in mathematics and physics is to identify a small set of generative principles from which large families of lawful structures may be derived. In the present work, that ambition is reframed in prime-recursive terms: construction is modeled by a typed tensor witness calculus, while stabilization is modeled by a contractive prime-weighted operator acting on a direct-sum tensor space. The central claim is not that arbitrary structure is generated \emph{ex nihilo}, but that admissible structure can be encoded, transformed, and stabilized within a common recursive formalism.

The framework is built from two layers. The first is the Prime-Encoded Tensor Constructor (\PETC), a provenance-aware tensor type system whose signatures live in the free abelian group over a set of atoms, and whose reshape legality is witnessed by explicit proof objects [file:12][file:13]. The second is the Universal Multiplicity Constant \(\Lambda_m\), understood here not as a primitive scalar but as a scalar reduction of a tensor-valued operator \(\Lambda_m^{\op}(t)\), with the scalar value used to enforce contractivity in a prime-recursive affine map.

This paper establishes the mathematically supported portion of that program. It proves reshape completeness in the \PETC\ layer, gives a precise operator/scalar separation for \(\Lambda_m\), proves Banach contractivity for the meta-recursive map under explicit assumptions, and derives a reachability corollary for admissible fixed points. It also records how commitment encodings and bytecode headers from the \(\mathrm{P}^2\mathrm{C}\) core can serve as a governance substrate for auditable tensor transformations [file:13].

Equally important is what this paper does \emph{not} claim. It does not establish a theory of consciousness, it does not replace proof theory, and it does not prove that all semantic or physical phenomena are exhaustively captured by the present framework. Those broader applications are treated as ongoing formalization rather than established theorem. This distinction is essential to preserving the technical integrity of the work.

\section{Mathematical Foundations of PETC}

\subsection{Signatures, axes, and tensor types}

Let \(A\) be a set of generators, called \emph{atoms}. Following the \(\mathrm{P}^2\mathrm{C}/\PETC\) core, define the signature group as the free abelian group on \(A\),
\[
\Sig := \Int^{(A)},
\]
that is, the finitely supported maps \(e : A \to \Int\) [file:12][file:13].

The positive cone is
\[
\Sig_+ := \Nat^{(A)} \subset \Sig,
\]
used when tracking positive multiplicities of axis atoms in reshape provenance [file:12].

A tensor axis may be written either in the unsigned form \(e \in \Sig\) or, in the variance-aware formalization, as a pair
\[
(\sigma,e) \in \{+1,-1\}\times \Sig_+,
\]
with \(\sigma\) indicating variance [file:12]. A tensor type or shape is then a finite list of such axes.

In the minimal witness-calculus presentation, a tensor type is simply a finite list of signatures,
\[
T = [e_1,\dots,e_n], \qquad e_i \in \Sig,
\]
with the production invariants that empty axes are removed, zero exponents are pruned, and scalar representation is the empty list [file:13].

\subsection{Global signature and structural invariants}

The central invariant is the \emph{global signature}
\[
\globalSig(T) := \sum_{i=1}^n e_i.
\]
In the variance-aware presentation this is written as the signed sum of axis contributions [file:12]. The conceptual role is the same in both versions: legal structural operations preserve, or predictably transform, global signature.

The primitive certified operations are:

\begin{itemize}[leftmargin=2em]
\item \textbf{Tensor product:} concatenation of axis lists, with global signature additive.
\item \textbf{Permutation:} reordering of axes, preserving global signature.
\item \textbf{Merge/Split:} refactoring of positive signatures into parts whose sum is conserved.
\item \textbf{Contraction:} cancellation of a pair of dual axes of opposite variance or opposite signature, preserving the appropriate invariant in the reduced tensor [file:12][file:13].
\end{itemize}

These operations are not merely operational conveniences. In \PETC, they are the generators of a witness language. A legal reshape is not ``size equality'' in the conventional numerical sense; it is a proof-carrying decomposition and recombination of signature mass [file:12][file:13].

\subsection{Weak and strong equivalence}

Two levels of type comparison are used in the \(\mathrm{P}^2\mathrm{C}\) core [file:13]:

\begin{itemize}[leftmargin=2em]
\item \textbf{Strong equality} preserves axis structure.
\item \textbf{Weak equivalence} is the reshape relation generated by \(\{\mathrm{Permute},\mathrm{Merge},\mathrm{Split}\}\).
\end{itemize}

We write strong equality as \(T_1 \strong T_2\) and weak equivalence as \(T_1 \weak T_2\). Then
\[
T_1 \strong T_2 \implies T_1 \weak T_2,
\]
but not conversely [file:13]. This distinction is useful in practice: strong equality is relevant for layout- and sharding-sensitive guarantees, while weak equivalence captures structural inter-convertibility under lawful reshape.

\subsection{Reshape completeness}

We now state the foundational theorem imported from the \(\mathrm{P}^2\mathrm{C}\) core.

\begin{theorem}[Reshape completeness]
Let \(T_1\) and \(T_2\) be finite axis-signature lists. Then
\[
T_1 \weak T_2 \iff \globalSig(T_1) = \globalSig(T_2).
\]
\end{theorem}

\begin{proof}
Each primitive operation \(\mathrm{Permute}\), \(\mathrm{Merge}\), and \(\mathrm{Split}\) preserves global signature, so \(T_1 \weak T_2\) implies \(\globalSig(T_1)=\globalSig(T_2)\). Conversely, if the global signatures agree, then each list may be merged to a single axis carrying the same total signature and then split into the target axes. Thus the relation generated by these primitives is exactly global-signature equality [file:13].
\end{proof}

This theorem is one of the key reasons \PETC\ is best interpreted as a universal constructor \emph{within its own domain}: any target tensor type weakly equivalent to a source type is constructible by an explicit witness synthesized from the primitive operations [file:13].

\section{The Universal Multiplicity Constant}

\subsection{The sealed two-layer operator}

We now introduce the stabilizing layer. Let \(P\) denote the set of primes and \(P_N\) the finite set of the first \(N\) primes. Define the tensor-valued operator
\[
\Lambda_m^{\op}(t)
=
M(\xi(p_i)) \circ M(\psi(p_i,t)).
\]
In the heuristic expansion used in the present program, this may be represented as
\[
\Lambda_m^{\op}(t)
=
\sum_{p_i\in P}
T_{ij}^{(p_i)}\, p_i^{\alpha_i}\bigl[\xi(p_i)+\psi(p_i,t)\bigr].
\]
Here \(\xi(p_i)\) is the static skeleton term and \(\psi(p_i,t)\) is the dynamic overlay.

The purpose of \(\Lambda_m^{\op}(t)\) is not to serve directly as the scalar contraction parameter. Rather, it is the full tensor-valued dynamic object from which a scalar contractivity bound is extracted.

\subsection{Scalar reduction}

We define the Universal Multiplicity Constant by spectral-radius reduction:
\[
\Lambda_m := \sup_t \spr\!\bigl(\Lambda_m^{\op}(t)\bigr).
\]
The contractive regime is the regime in which
\[
0 < \Lambda_m < 1.
\]

\begin{remark}
This separation resolves a frequent ambiguity. \(\Lambda_m^{\op}(t)\) is an operator-valued quantity encoding full tensorial dynamics, whereas \(\Lambda_m\) is the scalar bound used in the contraction parameter. The former retains structural richness; the latter enforces global contractivity.
\end{remark}

\subsection{The golden-ratio skeleton and its status}

A recurring choice in the broader Multiplicity program is
\[
\xi(p_i)=\log_{\Phi} p_i,
\]
where \(\Phi\) is the golden ratio. At present this paper treats that choice as a designated structural ansatz rather than as a completed theorem. It is retained because it is part of the intended geometry of the construction, but its exact orthogonality or projection-theoretic justification remains an open problem.

\begin{remark}
Accordingly, all results in this paper that depend on contractivity are formulated in terms of the scalar bound \(0<\Lambda_m<1\), not in terms of any unproved property specific to \(\xi(p_i)=\log_{\Phi}p_i\).
\end{remark}
\newpage
\section{The Prime-Recursive Meta-Map}

\subsection{Tensor space and prime decomposition}

Let
\[
\Hilb = \bigoplus_{p\in P}\Hilb_p
\]
be a prime-indexed direct-sum Hilbert space, and let \(\Tensor^{(m,n)}\) denote the Banach space of rank-\((m,n)\) tensors equipped with an operator norm induced by \(\Hilb\). For each prime \(p_i\), let \(\pi_{p_i}\) denote the orthogonal projection onto the \(p_i\)-component. Then
\[
\|\pi_{p_i}\|=1, \qquad \|T_{p_i}\|=\|\pi_{p_i}T\|\le \|T\|.
\]
This orthogonal decomposition is the analytic device used to control the prime-weighted sum.

The type-theoretic content of each \(T_{p_i}^{(m,n)}\) is governed by \PETC: its axes carry signature data in \(\Sig\), and its legal reshapes and contractions are constrained by the witness algebra and global-signature invariants [file:12][file:13].

\subsection{Definition of the affine map}

Fix \(\alpha < -1\), a finite prime truncation \(P_N\), and a driving term \(F^{(m,n)}\in \Tensor^{(m,n)}\). Define
\[
M(T)
=
A(T)+F^{(m,n)},
\]
where
\[
A(T)
=
\sum_{p_i\in P_N}\Lambda_m\, p_i^\alpha\, T_{p_i}^{(m,n)}.
\]
The associated scalar contraction parameter is
\[
k:=\sum_{p_i\in P_N}\Lambda_m\, p_i^\alpha.
\]
We work in the regime \(k<1\).

\begin{remark}
The condition \(\alpha<-1\) ensures convergence of the corresponding prime-weighted series and suppresses high-prime contributions. In this sense smaller primes dominate the recursive weighting, while finite truncation \(P_N\) makes the map simulation-ready.
\end{remark}

\subsection{Banach contractivity}

\begin{theorem}[Banach contractivity of the meta-recursive map]
Let \(\Tensor^{(m,n)}\) be the Banach space described above, and assume \(k<1\). Then the affine map
\[
M(T)=A(T)+F^{(m,n)}
\]
is a strict contraction with Lipschitz constant at most \(k\). Consequently, \(M\) has a unique fixed point \(T_\infty^{(m,n)}\in \Tensor^{(m,n)}\), and for every initial tensor \(T_0\),
\[
\|T_t-T_\infty\|\le k^t\|T_0-T_\infty\| \to 0
\qquad\text{as } t\to\infty,
\]
where \(T_{t+1}=M(T_t)\).
\end{theorem}

\begin{proof}
For any \(T,S\in \Tensor^{(m,n)}\),
\[
M(T)-M(S)=A(T)-A(S)
=
\sum_{p_i\in P_N}\Lambda_m\, p_i^\alpha\,(T_{p_i}-S_{p_i}).
\]
By the triangle inequality and the projection bound \(\|T_{p_i}-S_{p_i}\|\le \|T-S\|\),
\[
\|M(T)-M(S)\|
\le
\sum_{p_i\in P_N}\Lambda_m\, p_i^\alpha\, \|T_{p_i}-S_{p_i}\|
\le
\left(\sum_{p_i\in P_N}\Lambda_m\, p_i^\alpha\right)\|T-S\|
=
k\|T-S\|.
\]
Since \(k<1\), \(M\) is a strict contraction. The Banach fixed-point theorem then gives existence and uniqueness of a fixed point \(T_\infty\), and geometric convergence of the Picard iterates follows from the standard contraction estimate.
\end{proof}

\subsection{Closed form and the decoupled case}

The fixed-point equation is
\[
T_\infty = A(T_\infty)+F^{(m,n)},
\]
hence
\[
(\Id-A)T_\infty = F^{(m,n)},
\qquad
T_\infty = (\Id-A)^{-1}F^{(m,n)}.
\]
Because \(\|A\|\le k<1\), the Neumann series converges:
\[
(\Id-A)^{-1} = \sum_{j=0}^{\infty}A^j.
\]

\begin{remark}
The scalar expression
\[
T_\infty = \frac{F^{(m,n)}}{1-k}
\]
holds only in the decoupled case where \(A\) acts as scalar multiplication by \(k\) on the relevant subspace, or equivalently when the prime-indexed components are dynamically diagonalized in a way that reduces \(A\) to \(k\Id\). In the general case, the correct fixed-point formula is the operator expression \(T_\infty=(\Id-A)^{-1}F^{(m,n)}\).
\end{remark}

\subsection{Reachability and universal stable contraction}

\begin{corollary}[Universal stable contractor]
Let \(T^* \in \Tensor^{(m,n)}\) be admissible and let \(\varepsilon>0\). In the decoupled case \(A=k\Id\), choose
\[
F^{(m,n)}=(1-k)T^*.
\]
Then the unique fixed point satisfies
\[
T_\infty^{(m,n)}=T^*,
\]
and therefore \(\|T_\infty^{(m,n)}-T^*\|<\varepsilon\). More generally, in the non-decoupled case one may choose
\[
F^{(m,n)}=(\Id-A)T^*
\]
to obtain \(T_\infty^{(m,n)}=T^*\).
\end{corollary}

\begin{proof}
In the decoupled case, the scalar formula applies directly:
\[
T_\infty = \frac{F^{(m,n)}}{1-k} = \frac{(1-k)T^*}{1-k}=T^*.
\]
In the general case, substituting \(F^{(m,n)}=(\Id-A)T^*\) into the operator formula gives
\[
T_\infty=(\Id-A)^{-1}(\Id-A)T^* = T^*.
\]
\end{proof}

The point of the corollary is not unconstrained generation. Its content is that the contraction architecture stabilizes any admissible target by suitable choice of driving term. Universality resides in the contractor property.

\section{PETC as Constructor, \(\Lambda_m\) as Contractor}

The formal relationship between the two layers may now be stated clearly.

\begin{itemize}[leftmargin=2em]
\item \textbf{\PETC\ as constructor:} \PETC\ encodes tensors with provenance-aware axis signatures and supplies explicit witnesses for lawful reshape, split, merge, permutation, and contraction [file:12][file:13].
\item \textbf{\(\Lambda_m\) as contractor:} the scalar reduction \(\Lambda_m\) enforces a contractive bound on prime-recursive iteration, guaranteeing existence, uniqueness, and convergence of fixed points.
\end{itemize}

This duality is the central philosophical and mathematical thesis of the paper. Construction without contraction produces unrestricted proliferation of candidate structures. Contraction without construction produces only stabilization of whatever structure is already given. Their combination yields a lawful recursive substrate in which candidate tensor structures are both generable and stabilizable.

\section{Governance Substrate: Commitments, Witnesses, and Auditability}

The \(\mathrm{P}^2\mathrm{C}\) core also supplies a production-oriented witness and commitment framework. In particular, it defines:

\begin{itemize}[leftmargin=2em]
\item weak commitments based on global signature canonicalization,
\item strong commitments based on axis-structure canonicalization,
\item BLAKE2b digests with personalization,
\item a witness bytecode header with versioning and bounds-checkable structure [file:13].
\end{itemize}

These elements do not by themselves prove any metaphysical claim. What they do provide is an auditable governance substrate for tensor transformations.

\begin{remark}
This is the precise sense in which a prime-weighted execution hashing program (\PWEH) can be grounded in the current framework. The commitment layer binds legal tensor transformation traces to stable encodings and verifiable bytecode formats. A full formal theorem for \PWEH\ remains future work, but the commitment machinery needed for such a theorem is already visible in the \(\mathrm{P}^2\mathrm{C}\) core [file:13].
\end{remark}

\section{Simulation Program: ZRSD and Enhanced Tunneling}

The present paper is primarily mathematical. Nonetheless, the broader research program aims at simulation-oriented realization, especially through a certified \ZRSD\ master equation and contractively controlled tunneling dynamics. At this stage the status should be stated carefully.

\subsection{Simulation-ready structure}

The contractive map \(M\), finite prime truncation \(P_N\), and explicit fixed-point theory make the framework simulation-ready in the technical sense: one can specify \(F^{(m,n)}\), choose \(\Lambda_m\) and \(\alpha\), and iterate until convergence subject to the contractive bound \(k<1\).

\subsection{What remains open}

The following remain implementation and validation tasks rather than completed results:

\begin{itemize}[leftmargin=2em]
\item a fully specified \ZRSD\ master equation embedded in the same tensor space;
\item enhanced tunneling simulations with explicit fidelity curves and certification telemetry;
\item reproducible fixtures and test harnesses that emit standard contractivity diagnostics for reference ensembles;
\item empirical justification, if any, for the particular choice \(\xi(p_i)=\log_\Phi p_i\).
\end{itemize}

The mathematics of contraction does not depend on those future simulations; the simulations are meant to instantiate and test the framework, not to substitute for the theorems already proved.

\section{Implications for Mathematical Foundations}

The framework established here suggests a disciplined version of a foundational thesis: tensor-structured mathematical objects may be treated as occupants of a prime-indexed recursive space in which lawful construction and lawful stabilization are formally separated and then recombined.

Three consequences are worth emphasizing.

First, prime indexing provides a canonical hierarchy of weighted components. Even when the prime product encoding is treated only as an optional codec, the prime-labeled organization supplies a distinctive recursive architecture [file:12][file:13].

Second, the move from unconstrained infinite summation to finite truncation plus contractive control replaces metaphysical breadth with mathematical manageability. The finite set \(P_N\) gives tractability, while the \(k<1\) criterion gives stability.

Third, proof-carrying tensor operations and fixed-point stability belong to different but compatible layers. \PETC\ governs admissible structure by witness; \(\Lambda_m\) governs admissible dynamics by contraction. The result is not a replacement for proof theory, but a recursive substrate in which type-correct tensor transformation and stable iteration coexist in a single formal architecture.
\newpage
\section{Conclusion}

We have assembled a unified prime-recursive framework with two principal components: a constructive \PETC\ layer based on axis signatures in a free abelian group, and a contractive \(\Lambda_m\) layer obtained by scalar reduction of a tensor-valued operator. Within that architecture we proved reshape completeness for \PETC, Banach contractivity for the prime-recursive affine map, and a reachability corollary for admissible fixed points [file:12][file:13].

The resulting picture is mathematically precise and intentionally bounded. \PETC\ is universal as a constructor of lawful tensor reshapes within its type-theoretic domain; \(\Lambda_m\) is universal as a contractor of admissible tensor dynamics within the contractive regime. Together they define a recursive substrate for tensor-structured mathematical existence.

The broader ambitions of the Multiplicity program -- including semantic, cognitive, and physical unification -- remain open research directions. The present paper should therefore be read as a technical foundation: not the end of the framework, but the point at which its core constructive and contractive laws become rigorously visible.

\appendix

\section{Core PETC / \(\mathrm{P}^2\mathrm{C}\) ingredients}

For convenience, we summarize the imported ingredients from the attached \(\PETC\) and \(\mathrm{P}^2\mathrm{C}\) materials [file:12][file:13]:

\begin{enumerate}[leftmargin=2em]
\item \textbf{Signatures:} finitely supported integer-valued maps on atoms.
\item \textbf{Global signature:} additive invariant of a tensor type.
\item \textbf{Witness primitives:} \(\mathrm{Permute}\), \(\mathrm{Merge}\), \(\mathrm{Split}\), \(\mathrm{Contract}\).
\item \textbf{Reshape completeness:} weak equivalence iff equal global signature.
\item \textbf{Canonical forms:} weak and strong commitment encodings.
\item \textbf{Bytecode headers:} stable runtime-verifiable witness framing.
\end{enumerate}

\section{Open technical problems}

\begin{enumerate}[leftmargin=2em]
\item Give a full derivation or replace the ansatz \(\xi(p_i)=\log_\Phi p_i\) with a structurally justified alternative.
\item Prove sufficient conditions for \(0<\Lambda_m<1\) directly from properties of \(\Lambda_m^{\op}(t)\).
\item Complete a full formal theorem for \PWEH\ as an integrity functional.
\item Build reproducible simulation fixtures and a contractivity test harness for standard prime-indexed ensembles.
\item Port the witness algebra and contraction theorem to a proof assistant for machine-checked verification.
\end{enumerate}
\newpage
\section{Proof of Reshape Completeness}

\begin{definition}[Signature group and tensor types]
Let \(A\) be a set of atoms. The signature group is the free abelian group
\[
\Sig := \mathbb{Z}^{(A)},
\]
the finitely supported maps \(e : A \to \mathbb{Z}\). A tensor type is a finite list
\[
T = [e_1,\dots,e_n], \qquad e_i \in \Sig.
\]
Its global signature is
\[
\Sigma(T) := \sum_{i=1}^n e_i.
\]
These are the ground-truth objects in the \(\mathrm{P}^2\mathrm{C}/\mathrm{PETC}\) witness calculus [file:13].
\end{definition}

\begin{definition}[Weak equivalence]
Let \(\sim\) be the equivalence relation on tensor types generated by the primitive structural operations
\[
\{\mathrm{Permute},\mathrm{Merge},\mathrm{Split}\}.
\]
That is, \(T_1 \sim T_2\) iff \(T_2\) can be obtained from \(T_1\) by a finite sequence of these operations [file:13].
\end{definition}

\begin{theorem}[Reshape completeness]
For tensor types \(T_1,T_2\),
\[
T_1 \sim T_2 \iff \Sigma(T_1)=\Sigma(T_2).
\]
\end{theorem}

\begin{proof}
We prove both directions.

\paragraph{(\(\Rightarrow\)) Invariance under primitive operations.}
It suffices to show that each generator of \(\sim\) preserves global signature.

\begin{enumerate}
\item \textbf{Permutation.} If \(T=[e_1,\dots,e_n]\) and \(\pi\) is a permutation of \(\{1,\dots,n\}\), then
\[
\Sigma([e_{\pi(1)},\dots,e_{\pi(n)}])
=
\sum_{i=1}^n e_{\pi(i)}
=
\sum_{i=1}^n e_i
=
\Sigma(T).
\]

\item \textbf{Merge.} Suppose \(T\) contains entries \(e_{i_1},\dots,e_{i_r}\), and these are replaced by their sum
\[
e_* := e_{i_1}+\cdots+e_{i_r}.
\]
Then the new tensor type \(T'\) satisfies
\[
\Sigma(T')
=
\Sigma(T)-\sum_{j=1}^r e_{i_j}+e_*
=
\Sigma(T).
\]

\item \textbf{Split.} Suppose an entry \(e\) is replaced by a list of nonempty parts \(f_1,\dots,f_s\) such that
\[
e = f_1+\cdots+f_s.
\]
Then the new tensor type \(T'\) satisfies
\[
\Sigma(T')
=
\Sigma(T)-e+\sum_{j=1}^s f_j
=
\Sigma(T).
\]
\end{enumerate}

Therefore every primitive preserves \(\Sigma\), and hence every finite composition preserves \(\Sigma\). Thus
\[
T_1 \sim T_2 \implies \Sigma(T_1)=\Sigma(T_2).
\]

\paragraph{(\(\Leftarrow\)) Synthesis from global-signature equality.}
Assume now that
\[
\Sigma(T_1)=\Sigma(T_2)=:g.
\]
Write
\[
T_1=[e_1,\dots,e_m], \qquad T_2=[f_1,\dots,f_n].
\]

First merge all entries of \(T_1\) into a single axis carrying signature \(g\). This is possible by repeated application of \(\mathrm{Merge}\), yielding
\[
T_1 \sim [g].
\]

Likewise, since \(g=f_1+\cdots+f_n\), repeated application of \(\mathrm{Split}\) yields
\[
[g] \sim [f_1,\dots,f_n]=T_2.
\]

By transitivity of \(\sim\),
\[
T_1 \sim [g] \sim T_2,
\]
hence \(T_1 \sim T_2\).

This proves both directions.
\end{proof}

\begin{remark}
This theorem justifies the complete synthesis strategy implemented in the \(\mathrm{P}^2\mathrm{C}\) core: merge all source axes to their global signature, then split to the target axes. In the reference implementation this corresponds to the canonical weak form and the strategy \texttt{MergeAll ; SplitToTarget} [file:13].
\end{remark}

\section{Operator-Norm Bound for the Prime-Recursive Linear Part}

\begin{definition}[Prime-indexed direct-sum Hilbert space]
Let
\[
\mathcal{H} = \bigoplus_{p\in P} \mathcal{H}_p
\]
be an orthogonal direct sum of Hilbert spaces indexed by primes. For each \(p_i\), let
\[
\pi_{p_i} : \mathcal{H} \to \mathcal{H}_{p_i}
\]
denote the orthogonal projection onto the \(p_i\)-summand. Then
\[
\|\pi_{p_i}\| = 1.
\]
\end{definition}

\begin{lemma}[Projection bound]
For every \(T \in \mathcal{H}\),
\[
\| \pi_{p_i} T \| \le \|T\|.
\]
\end{lemma}

\begin{proof}
Since \(\pi_{p_i}\) is an orthogonal projection, it is contractive:
\[
\|\pi_{p_i}T\|^2
=
\langle \pi_{p_i}T,\pi_{p_i}T\rangle
\le
\langle T,T\rangle
=
\|T\|^2.
\]
Taking square roots gives the result.
\end{proof}

\begin{definition}[Prime-recursive linear part]
Let \(P_N\) be the first \(N\) primes, let \(0<\Lambda_m<1\), and let \(\alpha< -1\). Define the linear operator
\[
A(T) := \sum_{p_i\in P_N} \Lambda_m\, p_i^\alpha\, \pi_{p_i}T.
\]
Set
\[
k := \sum_{p_i\in P_N} \Lambda_m\, p_i^\alpha.
\]
\end{definition}

\begin{proposition}[Operator norm bound]
The linear operator \(A\) satisfies
\[
\|A\| \le k.
\]
In particular, if \(k<1\), then \(A\) is a strict contraction on the ambient Banach space.
\end{proposition}

\begin{proof}
Let \(T\) be arbitrary. Then by the triangle inequality,
\[
\|A(T)\|
=
\left\|
\sum_{p_i\in P_N}\Lambda_m\, p_i^\alpha\, \pi_{p_i}T
\right\|
\le
\sum_{p_i\in P_N}\Lambda_m\, p_i^\alpha\, \|\pi_{p_i}T\|.
\]
Applying the projection bound from the previous lemma,
\[
\|A(T)\|
\le
\sum_{p_i\in P_N}\Lambda_m\, p_i^\alpha\, \|T\|
=
k\|T\|.
\]
Taking the supremum over all \(T\neq 0\) gives
\[
\|A\|\le k.
\]
If \(k<1\), then \(\|A\|<1\), so \(A\) is a strict contraction.
\end{proof}

\begin{remark}
This is the precise norm estimate used in the Banach fixed-point argument for the affine map
\[
M(T)=A(T)+F.
\]
The proof depends on the orthogonal direct-sum decomposition and the fact that each \(\pi_{p_i}\) has norm \(1\).
\end{remark}

\section{Banach Fixed-Point Theorem for the Meta-Recursive Map}

\begin{theorem}[Contractivity of the affine map]
Let \(\mathcal{T}^{(m,n)}\) be a Banach space continuously embedded in the prime-indexed direct-sum Hilbert space, and define
\[
M(T)=A(T)+F^{(m,n)},
\]
where
\[
A(T)=\sum_{p_i\in P_N}\Lambda_m\, p_i^\alpha\, T_{p_i}^{(m,n)},
\qquad
k=\sum_{p_i\in P_N}\Lambda_m\, p_i^\alpha<1.
\]
Then for all \(T,S\in\mathcal{T}^{(m,n)}\),
\[
\|M(T)-M(S)\| \le k\|T-S\|.
\]
Hence \(M\) is a strict contraction and has a unique fixed point \(T_\infty^{(m,n)}\).
\end{theorem}

\begin{proof}
For any \(T,S\in \mathcal{T}^{(m,n)}\),
\[
M(T)-M(S)=A(T)-A(S).
\]
Therefore
\[
\|M(T)-M(S)\|
=
\left\|
\sum_{p_i\in P_N}\Lambda_m\, p_i^\alpha\, (T_{p_i}^{(m,n)}-S_{p_i}^{(m,n)})
\right\|.
\]
By the triangle inequality,
\[
\|M(T)-M(S)\|
\le
\sum_{p_i\in P_N}\Lambda_m\, p_i^\alpha\, \|T_{p_i}^{(m,n)}-S_{p_i}^{(m,n)}\|.
\]
By the orthogonal projection bound,
\[
\|T_{p_i}^{(m,n)}-S_{p_i}^{(m,n)}\|
\le
\|T-S\|.
\]
Hence
\[
\|M(T)-M(S)\|
\le
\left(\sum_{p_i\in P_N}\Lambda_m\, p_i^\alpha\right)\|T-S\|
=
k\|T-S\|.
\]
Since \(k<1\), \(M\) is a strict contraction. The Banach fixed-point theorem now gives existence and uniqueness of a fixed point \(T_\infty^{(m,n)}\), together with convergence of Picard iterates
\[
T_{t+1}=M(T_t).
\]
\end{proof}

\begin{corollary}[Geometric convergence]
For every initial tensor \(T_0\in\mathcal{T}^{(m,n)}\), the iterates \(T_t\) satisfy
\[
\|T_t-T_\infty\| \le k^t \|T_0-T_\infty\|,
\]
and therefore
\[
\|T_t-T_\infty\| \to 0
\qquad \text{as } t\to\infty.
\]
\end{corollary}

\begin{proof}
Using the contraction inequality and the fixed-point property \(M(T_\infty)=T_\infty\),
\[
\|T_{t+1}-T_\infty\|
=
\|M(T_t)-M(T_\infty)\|
\le
k\|T_t-T_\infty\|.
\]
Iterating this estimate yields
\[
\|T_t-T_\infty\|
\le
k^t\|T_0-T_\infty\|.
\]
Since \(0\le k<1\), the right-hand side tends to \(0\).
\end{proof}

\section{Closed-Form Fixed Point and the Decoupled Case}

\begin{proposition}[General fixed-point formula]
Let \(M(T)=A(T)+F\) with \(\|A\|<1\). Then the unique fixed point is
\[
T_\infty = (\mathrm{Id}-A)^{-1}F,
\]
where
\[
(\mathrm{Id}-A)^{-1} = \sum_{j=0}^\infty A^j
\]
in operator norm.
\end{proposition}

\begin{proof}
The fixed-point equation is
\[
T_\infty = A(T_\infty)+F.
\]
Rearranging,
\[
(\mathrm{Id}-A)T_\infty = F.
\]
Since \(\|A\|<1\), the Neumann series converges in operator norm:
\[
(\mathrm{Id}-A)^{-1} = \sum_{j=0}^{\infty}A^j.
\]
Applying this inverse gives
\[
T_\infty = (\mathrm{Id}-A)^{-1}F.
\]
\end{proof}

\begin{proposition}[Decoupled scalar case]
Assume in addition that on the relevant invariant subspace one has
\[
A = k\,\mathrm{Id}
\qquad\text{with } 0\le k<1.
\]
Then
\[
T_\infty = \frac{F}{1-k}.
\]
\end{proposition}

\begin{proof}
If \(A=k\,\mathrm{Id}\), then
\[
(\mathrm{Id}-A) = (1-k)\mathrm{Id},
\]
so
\[
(\mathrm{Id}-A)^{-1} = \frac{1}{1-k}\,\mathrm{Id}.
\]
Substituting into the general formula yields
\[
T_\infty
=
(\mathrm{Id}-A)^{-1}F
=
\frac{1}{1-k}F.
\]
\end{proof}

\begin{remark}
This appendix isolates the exact condition under which the scalar closed form is valid. In the general prime-recursive setting, the correct fixed-point formula is the operator expression \(T_\infty=(\mathrm{Id}-A)^{-1}F\); the scalar formula \(F/(1-k)\) is a special decoupled case, not a universal identity.
\end{remark}

\section{Reachability of Admissible Targets}

\begin{corollary}[Reachability in the general operator case]
Let \(T^*\in \mathcal{T}^{(m,n)}\) be any admissible target. Assume \(\|A\|<1\). Then choosing
\[
F := (\mathrm{Id}-A)T^*
\]
forces \(T^*\) to be the unique fixed point of \(M(T)=A(T)+F\).
\end{corollary}

\begin{proof}
Substitute the chosen forcing term into the fixed-point formula:
\[
T_\infty
=
(\mathrm{Id}-A)^{-1}F
=
(\mathrm{Id}-A)^{-1}(\mathrm{Id}-A)T^*
=
T^*.
\]
Thus \(T^*\) is the unique fixed point.
\end{proof}

\begin{corollary}[Reachability in the decoupled scalar case]
If \(A=k\,\mathrm{Id}\) with \(0\le k<1\), then choosing
\[
F=(1-k)T^*
\]
forces
\[
T_\infty=T^*.
\]
\end{corollary}

\begin{proof}
Immediate from the scalar closed-form formula of Appendix D.
\end{proof}

\section{PETC Witness Legality and Canonical Weak Form}

\begin{proposition}[Canonical weak form]
For a tensor type \(T=[e_1,\dots,e_n]\), define its canonical weak form by collecting the nonzero entries of the global signature:
\[
\mathrm{cw}(T)
=
[(a_1,\Sigma(T)(a_1)),\dots,(a_r,\Sigma(T)(a_r))]
\]
over the support of \(\Sigma(T)\), ordered by atom index. Then
\[
T_1 \sim T_2 \iff \mathrm{cw}(T_1)=\mathrm{cw}(T_2).
\]
\end{proposition}

\begin{proof}
By reshape completeness,
\[
T_1\sim T_2 \iff \Sigma(T_1)=\Sigma(T_2).
\]
But \(\mathrm{cw}(T)\) is just a canonical encoding of \(\Sigma(T)\) as a pruned ordered list. Hence equality of canonical weak forms is equivalent to equality of global signatures.
\end{proof}

\begin{remark}
This is the mathematical justification for the implementation rule
\[
\texttt{reshape\_equivalent(src,tgt) := canonical\_weak(src) == canonical\_weak(tgt)}
\]
used in the \(\mathrm{P}^2\mathrm{C}\) production core [file:13].
\end{remark}

\nocite{*}
\bibliographystyle{plain}
\bibliography{references}

\end{document}

